\sectionheading{5.34 Python Data Type}

Program-এ ব্যবহৃত কোনো Value কী ধরনের এবং সেই Value-এর
উপর কী ধরনের Operation করা যাবে—তা
\textbf{Data Type} দ্বারা বোঝানো হয়।

%================================================
% Definition + first example
%================================================

\booksubheading{Data Type কী?}

\begin{definitionbox}
কোনো Data বা Value-এর
\textbf{ধরন, বৈশিষ্ট্য এবং তার উপর প্রযোজ্য Operation}
নির্দেশকারী শ্রেণিকে \textbf{Data Type} বলে।
\end{definitionbox}

\noindent
\begin{minipage}[t]{0.48\linewidth}

\textbf{উদাহরণ—}

\begin{codebox}
\begin{english}
\ttfamily
age = 18\\
height = 5.7\\
name = "Rahim"
\end{english}
\end{codebox}

\end{minipage}
\hfill
\begin{minipage}[t]{0.50\linewidth}

\begin{tabularx}
{\linewidth}
{
|>{\RaggedRight\arraybackslash}p{0.28\linewidth}
|>{\RaggedRight\arraybackslash}p{0.25\linewidth}
|Y|
}
\hline
\rowcolor{softBlue}

\textbf{Variable}
&
\textbf{Value}
&
\textbf{Type}
\\
\hline

\textenglish{\ttfamily age}
&
\textenglish{\ttfamily 18}
&
\textenglish{\ttfamily int}
\\
\hline

\textenglish{\ttfamily height}
&
\textenglish{\ttfamily 5.7}
&
\textenglish{\ttfamily float}
\\
\hline

\textenglish{\ttfamily name}
&
\textenglish{\ttfamily "Rahim"}
&
\textenglish{\ttfamily str}
\\
\hline
\end{tabularx}

\end{minipage}

\begin{rememberbox}
\begin{center}
\textbf{Value}
\(\longrightarrow\)
\textbf{কী ধরনের?}
\(\longrightarrow\)
\textbf{Data Type}
\end{center}
\end{rememberbox}

%================================================
% Main classification
%================================================

\booksubheading{Python-এর প্রধান Built-in Data Type}

\begin{tabularx}
{\linewidth}
{
|>{\bfseries\RaggedRight\arraybackslash}p{0.25\linewidth}
|Y|
}
\hline
\rowcolor{softBlue}

\multicolumn{1}{|c|}{\textbf{শ্রেণি}}
&
\multicolumn{1}{c|}{\textbf{Data Type}}
\\
\hline

Numeric
&
\textenglish{\ttfamily int, float, complex}
\\
\hline

Boolean
&
\textenglish{\ttfamily bool}
\\
\hline

Text
&
\textenglish{\ttfamily str}
\\
\hline

Sequence
&
\textenglish{\ttfamily list, tuple}
\\
\hline

Set
&
\textenglish{\ttfamily set}
\\
\hline

Mapping
&
\textenglish{\ttfamily dict}
\\
\hline

None
&
\textenglish{\ttfamily NoneType}
\\
\hline
\end{tabularx}

%================================================
% type() compact
%================================================

\booksubheading{\textenglish{\ttfamily type()} Function}

কোনো Value বা Variable-এর Data Type জানার জন্য
\textenglish{\ttfamily type()}
Function ব্যবহার করা যায়।

\noindent
\begin{minipage}[t]{0.49\linewidth}

\begin{examplebox}{Integer Type}

\begin{codebox}
\begin{english}
\ttfamily
x = 10\\
print(type(x))
\end{english}
\end{codebox}

Output:
\textenglish{\ttfamily \textless class 'int'\textgreater}

\end{examplebox}

\end{minipage}
\hfill
\begin{minipage}[t]{0.49\linewidth}

\begin{examplebox}{String Type}

\begin{codebox}
\begin{english}
\ttfamily
name = "Mina"\\
print(type(name))
\end{english}
\end{codebox}

Output:
\textenglish{\ttfamily \textless class 'str'\textgreater}

\end{examplebox}

\end{minipage}

\begin{rememberbox}
\textbf{\color{deepOrange}MCQ তথ্য:}
\quad
\textenglish{\ttfamily type()}
\(\longrightarrow\)
Data Type নির্ণয় করে।
\end{rememberbox}

%================================================
% INT + FLOAT side by side
%================================================

\chFiveInput{topics/34_python_data_types/python_extensions.tex}

\booksubheading{Numeric Data Type}

\noindent
\begin{minipage}[t]{0.49\linewidth}

\begin{examplebox}{১. Integer — \textenglish{\ttfamily int}}

পূর্ণসংখ্যার জন্য
\textenglish{\ttfamily int}
ব্যবহৃত হয়।

\begin{codebox}
\begin{english}
\ttfamily
a = 10\\
b = -25\\
c = 0
\end{english}
\end{codebox}

এগুলো সব
\textenglish{\ttfamily int}।

\begin{codebox}
\begin{english}
\ttfamily
print(type(a))
\end{english}
\end{codebox}

Output:
\textenglish{\ttfamily \textless class 'int'\textgreater}

\medskip

\textbf{Size/Range:}

Python
\textenglish{\ttfamily int}
সাধারণ fixed 32-bit Range-এ সীমাবদ্ধ নয়।

\begin{codebox}
\begin{english}
\ttfamily
x = 999999999999999999999999999999
\end{english}
\end{codebox}

এটিও
\textenglish{\ttfamily int}
হতে পারে।

\textbf{ব্যবহারিক সীমা}
\(\longrightarrow\)
Available Memory।

\end{examplebox}

\end{minipage}
\hfill
\begin{minipage}[t]{0.49\linewidth}

\begin{examplebox}{২. Floating-point — \textenglish{\ttfamily float}}

দশমিক বা Floating-point সংখ্যার জন্য
\textenglish{\ttfamily float}
ব্যবহৃত হয়।

\begin{codebox}
\begin{english}
\ttfamily
x = 10.5\\
y = -3.75\\
z = 2.0
\end{english}
\end{codebox}

এগুলো সব
\textenglish{\ttfamily float}।

\begin{codebox}
\begin{english}
\ttfamily
print(type(x))
\end{english}
\end{codebox}

Output:
\textenglish{\ttfamily \textless class 'float'\textgreater}

\medskip

\textbf{Size/Range:}

সাধারণ CPython-এ
\textenglish{\ttfamily float}
সাধারণত 64-bit double-precision floating-point representation।

প্রায় সর্বোচ্চ magnitude:
\[
\pm 1.8\times10^{308}
\]

খুব ছোট Positive Normal Value:
\[
2.2\times10^{-308}
\]

\end{examplebox}

\end{minipage}

\begin{rememberbox}
\textbf{\color{deepOrange}সতর্কতা:}
Python
\textenglish{\ttfamily int}
-এর fixed 4-byte/32-bit range নেই।
আর
\textenglish{\ttfamily float}
-এর 64-bit বলতে floating-point representation বোঝানো হয়,
পুরো Python Object-এর Memory Size নয়।
\end{rememberbox}

%================================================
% Complex + Bool
%================================================

\noindent
\begin{minipage}[t]{0.49\linewidth}

\begin{examplebox}{৩. Complex — \textenglish{\ttfamily complex}}

Complex Number-এর জন্য
\textenglish{\ttfamily complex}
ব্যবহৃত হয়।

\begin{codebox}
\begin{english}
\ttfamily
z = 3 + 4j\\
print(type(z))
\end{english}
\end{codebox}

এখানে—

\begin{itemize}
\item Real Part
\(\longrightarrow\)
\textenglish{\ttfamily 3}

\item Imaginary Part
\(\longrightarrow\)
\textenglish{\ttfamily 4j}
\end{itemize}

Output:
\textenglish{\ttfamily \textless class 'complex'\textgreater}

\medskip

\textbf{MCQ:}
Python Complex Number-এর Imaginary অংশে
\textenglish{\ttfamily j}
ব্যবহৃত হয়।

\end{examplebox}

\end{minipage}
\hfill
\begin{minipage}[t]{0.49\linewidth}

\begin{examplebox}{৪. Boolean — \textenglish{\ttfamily bool}}

Logical সত্য-মিথ্যা Value-এর জন্য
\textenglish{\ttfamily bool}
ব্যবহৃত হয়।

\begin{center}
\textenglish{\ttfamily True}
\qquad
\textenglish{\ttfamily False}
\end{center}

\begin{codebox}
\begin{english}
\ttfamily
is\_student = True\\
print(type(is\_student))
\end{english}
\end{codebox}

Output:
\textenglish{\ttfamily \textless class 'bool'\textgreater}

\medskip

\textbf{মনে রাখবে—}

\begin{center}
\textenglish{\ttfamily bool}
\(\longrightarrow\)
\textenglish{\ttfamily True, False}
\end{center}

\textenglish{\ttfamily True}
ও
\textenglish{\ttfamily False}
-এর প্রথম Letter Capital।

\end{examplebox}

\end{minipage}

%================================================
% String + NoneType
%================================================

\booksubheading{Text ও বিশেষ Data Type}

\noindent
\begin{minipage}[t]{0.49\linewidth}

\begin{examplebox}{৫. String — \textenglish{\ttfamily str}}

Text বা Character-এর ধারাবাহিকতার জন্য
\textenglish{\ttfamily str}
ব্যবহৃত হয়।

\begin{codebox}
\begin{english}
\ttfamily
name = "Rahim"\\
country = 'Bangladesh'
\end{english}
\end{codebox}

String-এ Single Quote বা Double Quote ব্যবহার করা যায়।

\begin{center}
\textenglish{\ttfamily 'Mina'}
\qquad অথবা \qquad
\textenglish{\ttfamily "Mina"}
\end{center}

\textbf{Size:}
Fixed Size নেই; Character-এর সংখ্যা ও Internal Representation
অনুযায়ী Memory Usage পরিবর্তিত হয়।

\end{examplebox}

\end{minipage}
\hfill
\begin{minipage}[t]{0.49\linewidth}

\begin{examplebox}{১০. NoneType}

কোনো অর্থপূর্ণ Value বর্তমানে নির্ধারিত নেই—এমন বিশেষ অবস্থা
প্রকাশ করতে
\textenglish{\ttfamily None}
ব্যবহৃত হয়।

\begin{codebox}
\begin{english}
\ttfamily
x = None\\
print(type(x))
\end{english}
\end{codebox}

Output:
\textenglish{\ttfamily \textless class 'NoneType'\textgreater}

\medskip

\begin{center}
\textenglish{\ttfamily None}
\(\neq\)
\textenglish{\ttfamily 0}
\qquad
\textenglish{\ttfamily None}
\(\neq\)
\textenglish{\ttfamily False}

\medskip

\textenglish{\ttfamily None}
\(\neq\)
\textenglish{\ttfamily ""}
\end{center}

\textenglish{\ttfamily None}
একটি আলাদা বিশেষ Value।

\end{examplebox}

\end{minipage}

%================================================
% LIST + TUPLE
%================================================

\booksubheading{Sequence Data Type}

\noindent
\begin{minipage}[t]{0.49\linewidth}

\begin{examplebox}{৬. List — \textenglish{\ttfamily list}}

একাধিক Value একটি Collection-এ রাখার জন্য List ব্যবহার করা যায়।

\begin{codebox}
\begin{english}
\ttfamily
numbers = [10, 20, 30]\\
data = [10, "Rahim", 5.5, True]
\end{english}
\end{codebox}

\textbf{বৈশিষ্ট্য—}

\begin{itemize}
\item Ordered
\item Mutable বা পরিবর্তনযোগ্য
\item Duplicate Value থাকতে পারে
\end{itemize}

\begin{codebox}
\begin{english}
\ttfamily
numbers[0] = 100
\end{english}
\end{codebox}

অর্থাৎ List-এর Element পরিবর্তন করা যায়।

\end{examplebox}

\end{minipage}
\hfill
\begin{minipage}[t]{0.49\linewidth}

\begin{examplebox}{৭. Tuple — \textenglish{\ttfamily tuple}}

Tuple-এও একাধিক Value রাখা যায়।

\begin{codebox}
\begin{english}
\ttfamily
numbers = (10, 20, 30)
\end{english}
\end{codebox}

\textbf{বৈশিষ্ট্য—}

\begin{itemize}
\item Ordered
\item Immutable বা অপরিবর্তনযোগ্য
\item Duplicate Value থাকতে পারে
\end{itemize}

\medskip

\begin{center}
\textbf{List}
\(\longrightarrow\)
Mutable
\qquad
\textbf{Tuple}
\(\longrightarrow\)
Immutable
\end{center}

\end{examplebox}

\end{minipage}

%================================================
% List vs Tuple table
%================================================

\begin{tabularx}
{\linewidth}
{
|>{\bfseries\RaggedRight\arraybackslash}p{0.22\linewidth}
|Y
|Y|
}
\hline
\rowcolor{softBlue}

\multicolumn{1}{|c|}{\textbf{বিষয়}}
&
\multicolumn{1}{c|}{\textbf{List}}
&
\multicolumn{1}{c|}{\textbf{Tuple}}
\\
\hline

Syntax
&
\textenglish{\ttfamily [ ]}
&
\textenglish{\ttfamily ( )}
\\
\hline

পরিবর্তন
&
Mutable
&
Immutable
\\
\hline

Ordered
&
হ্যাঁ
&
হ্যাঁ
\\
\hline

Duplicate
&
থাকতে পারে
&
থাকতে পারে
\\
\hline

উদাহরণ
&
\textenglish{\ttfamily [1,2,3]}
&
\textenglish{\ttfamily (1,2,3)}
\\
\hline
\end{tabularx}

%================================================
% SET + DICT — Space Optimized
%================================================

\booksubheading{Set ও Mapping Data Type}

\noindent
\begin{minipage}[t]{0.49\linewidth}

\begin{examplebox}{৮. Set — \textenglish{\ttfamily set}}

একাধিক \textbf{Unique Value}-এর Collection তৈরিতে
\textenglish{\ttfamily set}
ব্যবহৃত হয়।

\begin{codebox}
\begin{english}
\ttfamily
a = \{10, 20, 30\}\\
b = \{10, 20, 20, 30\}
\end{english}
\end{codebox}

দ্বিতীয় Set-এ Duplicate
\textenglish{\ttfamily 20}
আলাদাভাবে সংরক্ষিত হয় না।

\begin{center}
\textenglish{\ttfamily \{10, 20, 30\}}
\end{center}

\textbf{বৈশিষ্ট্য:}\quad
Unique
\(\;\bullet\;\)
Mutable
\(\;\bullet\;\)
সাধারণ Indexed Sequence নয়

\end{examplebox}

\end{minipage}
\hfill
\begin{minipage}[t]{0.49\linewidth}

\begin{examplebox}{৯. Dictionary — \textenglish{\ttfamily dict}}

Data-কে
\textbf{Key : Value}
Pair আকারে সংরক্ষণ করতে
\textenglish{\ttfamily dict}
ব্যবহৃত হয়।

\begin{codebox}
\begin{english}
\ttfamily
student = \{"name": "Rahim", "age": 18\}
\end{english}
\end{codebox}

\begin{center}

\textenglish{\ttfamily "name"}
\(\longrightarrow\)
Key
\qquad
\textenglish{\ttfamily "Rahim"}
\(\longrightarrow\)
Value

\medskip

\textenglish{\ttfamily "age"}
\(\longrightarrow\)
Key
\qquad
\textenglish{\ttfamily 18}
\(\longrightarrow\)
Value

\end{center}

\textbf{বৈশিষ্ট্য:}\quad
Key--Value Pair
\(\;\bullet\;\)
Mutable

\end{examplebox}

\end{minipage}

\begin{rememberbox}

\textbf{\color{deepOrange}MCQ/CQ-এর জন্য মনে রাখবে}

\begin{center}

\textbf{Set}
\(\longrightarrow\)
Unique Value-এর Collection

\qquad

\textbf{Dictionary}
\(\longrightarrow\)
Key : Value Pair

\medskip

\textbf{Set ও Dictionary উভয়ই}
\(\longrightarrow\)
\textbf{Mutable}

\end{center}

\end{rememberbox}

%================================================
% Size / Range Table
%================================================

\booksubheading{Python Data Type-এর Size ও Range}

Python-এ C/C++-এর মতো সব Data Type-এর জন্য নির্দিষ্ট
Fixed Byte Size ও Fixed Range দেওয়া সঠিক নয়।

\begin{tabularx}
{\linewidth}
{
|>{\bfseries\RaggedRight\arraybackslash}p{0.20\linewidth}
|Y|
}
\hline
\rowcolor{softBlue}

\multicolumn{1}{|c|}{\textbf{Data Type}}
&
\multicolumn{1}{c|}{\textbf{Size / Range-এর ধারণা}}
\\
\hline

\textenglish{\ttfamily int}
&
Fixed Range নেই; প্রয়োজন অনুযায়ী বড় হতে পারে।
ব্যবহারিক সীমা available memory-এর উপর নির্ভর করে।
\\
\hline

\textenglish{\ttfamily float}
&
সাধারণ CPython-এ সাধারণত 64-bit floating-point representation;
প্রায় \(1.8\times10^{308}\) magnitude পর্যন্ত।
\\
\hline

\textenglish{\ttfamily complex}
&
Real ও Imaginary—দুটি floating-point অংশ নিয়ে গঠিত।
\\
\hline

\textenglish{\ttfamily bool}
&
শুধু
\textenglish{\ttfamily True}
ও
\textenglish{\ttfamily False}।
\\
\hline

\textenglish{\ttfamily str}
&
Fixed Size নেই; Content অনুযায়ী Memory Usage পরিবর্তিত হয়।
\\
\hline

\textenglish{\ttfamily list}
&
Fixed Size নেই; Element বাড়লে Memory Usage বাড়ে।
\\
\hline

\textenglish{\ttfamily tuple}
&
Fixed Size নেই; Element-এর সংখ্যা ও Content-এর উপর নির্ভর করে।
\\
\hline

\textenglish{\ttfamily set}
&
Fixed Size নেই; Element-এর সংখ্যা অনুযায়ী পরিবর্তিত হয়।
\\
\hline

\textenglish{\ttfamily dict}
&
Fixed Size নেই; Key-Value Pair অনুযায়ী Memory Usage পরিবর্তিত হয়।
\\
\hline

\textenglish{\ttfamily NoneType}
&
বিশেষ
\textenglish{\ttfamily None}
Value নির্দেশ করে।
\\
\hline
\end{tabularx}

\begin{rememberbox}

\textbf{\color{deepOrange}খুব গুরুত্বপূর্ণ}

Python-এর জন্য এভাবে লিখবে না—

\begin{center}
\textenglish{\ttfamily int = 4 byte}
\qquad
\textenglish{\ttfamily Range = -2\textasciicircum31 to 2\textasciicircum31 - 1}
\end{center}

এটি Python
\textenglish{\ttfamily int}
-এর সাধারণ নিয়ম নয়।

\begin{center}

\textbf{C/C++}
\(\longrightarrow\)
অনেক Numeric Type-এর নির্দিষ্ট Size/Range

\medskip

\textbf{Python}
\(\longrightarrow\)
অনেক ক্ষেত্রে Size Dynamic বা Implementation-dependent

\end{center}

\end{rememberbox}

%================================================
% Mutable Immutable
%================================================

\booksubheading{Mutable ও Immutable Data Type}

Data তৈরি হওয়ার পর তার Content পরিবর্তন করা যায় কি না—
তার ভিত্তিতে Data Type-কে Mutable ও Immutable হিসেবে বোঝা যায়।

\noindent
\begin{minipage}[t]{0.49\linewidth}

\begin{definitionbox}

\textbf{\color{darkTeal}Mutable}

Object তৈরি হওয়ার পর তার Content পরিবর্তন করা যায়।

\begin{center}
\textenglish{\ttfamily list}
\qquad
\textenglish{\ttfamily set}
\qquad
\textenglish{\ttfamily dict}
\end{center}

\end{definitionbox}

\end{minipage}
\hfill
\begin{minipage}[t]{0.49\linewidth}

\begin{definitionbox}

\textbf{\color{darkTeal}Immutable}

Object তৈরি হওয়ার পর তার নিজস্ব Value/Content পরিবর্তন করা যায় না।

\begin{center}
\textenglish{\ttfamily int, float, complex, bool, str, tuple, NoneType}
\end{center}

\end{definitionbox}

\end{minipage}

\begin{rememberbox}
\begin{center}

\textbf{List, Set, Dictionary}
\(\longrightarrow\)
\textbf{Mutable}

\qquad

\textbf{Number, String, Tuple}
\(\longrightarrow\)
\textbf{Immutable}

\end{center}
\end{rememberbox}

%================================================
% Summary table
%================================================

\booksubheading{Python Data Type এক নজরে}

\begin{tabularx}
{\linewidth}
{
|>{\bfseries\RaggedRight\arraybackslash}p{0.17\linewidth}
|>{\RaggedRight\arraybackslash}p{0.25\linewidth}
|Y|
}
\hline
\rowcolor{softBlue}

\multicolumn{1}{|c|}{\textbf{Type}}
&
\multicolumn{1}{c|}{\textbf{উদাহরণ}}
&
\multicolumn{1}{c|}{\textbf{মূল কাজ}}
\\
\hline

\textenglish{\ttfamily int}
&
\textenglish{\ttfamily 10}
&
পূর্ণসংখ্যা
\\
\hline

\textenglish{\ttfamily float}
&
\textenglish{\ttfamily 10.5}
&
দশমিক/Floating-point সংখ্যা
\\
\hline

\textenglish{\ttfamily complex}
&
\textenglish{\ttfamily 3+4j}
&
Complex Number
\\
\hline

\textenglish{\ttfamily bool}
&
\textenglish{\ttfamily True}
&
Logical Value
\\
\hline

\textenglish{\ttfamily str}
&
\textenglish{\ttfamily "Rahim"}
&
Text
\\
\hline

\textenglish{\ttfamily list}
&
\textenglish{\ttfamily [1,2,3]}
&
Mutable Collection
\\
\hline

\textenglish{\ttfamily tuple}
&
\textenglish{\ttfamily (1,2,3)}
&
Immutable Collection
\\
\hline

\textenglish{\ttfamily set}
&
\textenglish{\ttfamily \{1,2,3\}}
&
Unique Element-এর Collection
\\
\hline

\textenglish{\ttfamily dict}
&
\textenglish{\ttfamily \{"a":1\}}
&
Key-Value Collection
\\
\hline

\textenglish{\ttfamily NoneType}
&
\textenglish{\ttfamily None}
&
বিশেষ No-value অবস্থা
\\
\hline

\end{tabularx}

%================================================
% Mixed program
%================================================

\booksubheading{একই Program-এ বিভিন্ন Data Type}

\begin{examplebox}{উদাহরণ}

\begin{codebox}
\begin{english}
\ttfamily
name = "Mina"\\
age = 18\\
height = 5.4\\
student = True\\
marks = [80, 85, 90]\\[0.5mm]
print(type(name))\\
print(type(age))\\
print(type(height))\\
print(type(student))\\
print(type(marks))
\end{english}
\end{codebox}

Data Type যথাক্রমে—

\begin{center}
\textenglish{\ttfamily str}
\(\longrightarrow\)
\textenglish{\ttfamily int}
\(\longrightarrow\)
\textenglish{\ttfamily float}
\(\longrightarrow\)
\textenglish{\ttfamily bool}
\(\longrightarrow\)
\textenglish{\ttfamily list}
\end{center}

\end{examplebox}

%================================================
% Common mistakes
%================================================

\booksubheading{সাধারণ ভুল}

\begin{compacttwocolpoints}

\item
Python
\textenglish{\ttfamily int}
-কে সবসময় 4 Byte মনে করা।

\item
Python
\textenglish{\ttfamily int}
-এর Fixed 32-bit Range আছে মনে করা।

\item
\textenglish{\ttfamily "10"}
ও
\textenglish{\ttfamily 10}
-কে একই Data Type মনে করা।

\item
\textenglish{\ttfamily 3 + 4j}
-এর
\textenglish{\ttfamily j}
-কে সাধারণ Variable মনে করা।

\item
\textenglish{\ttfamily True}
-এর পরিবর্তে
\textenglish{\ttfamily true}
লেখা।

\item List ও Tuple-কে একই ধরনের পরিবর্তনযোগ্য Collection মনে করা।

\item Set-এ Duplicate Element আলাদাভাবে থাকবে মনে করা।

\item Dictionary-কে শুধু Value-এর Collection মনে করা।

\item
\textenglish{\ttfamily None}
-কে
\textenglish{\ttfamily 0}
মনে করা।

\item Mutable ও Immutable-এর পার্থক্য না বোঝা।

\end{compacttwocolpoints}

%================================================
% Final Reminder
%================================================

\begin{rememberbox}

\textbf{\color{deepOrange}মনে রাখবে}

\begin{tabularx}
{\linewidth}
{
>{\bfseries\RaggedRight\arraybackslash}p{0.23\linewidth}
Y
>{\bfseries\RaggedRight\arraybackslash}p{0.23\linewidth}
Y
}

\textenglish{\ttfamily int}
&
Integer
&
\textenglish{\ttfamily float}
&
Floating-point
\\

\textenglish{\ttfamily complex}
&
\textenglish{\ttfamily a+bj}
&
\textenglish{\ttfamily bool}
&
\textenglish{\ttfamily True/False}
\\

\textenglish{\ttfamily str}
&
Text
&
\textenglish{\ttfamily list}
&
Mutable Collection
\\

\textenglish{\ttfamily tuple}
&
Immutable Collection
&
\textenglish{\ttfamily set}
&
Unique Collection
\\

\textenglish{\ttfamily dict}
&
Key : Value
&
\textenglish{\ttfamily None}
&
NoneType
\\

\textenglish{\ttfamily type()}
&
Data Type নির্ণয়
&
Python \textenglish{\ttfamily int}
&
Fixed 4-byte Range নয়
\\

\end{tabularx}

\medskip

\begin{center}

\textbf{Number}
\(\longrightarrow\)
\textenglish{\ttfamily int, float, complex}
\qquad
\textbf{Text}
\(\longrightarrow\)
\textenglish{\ttfamily str}

\medskip

\textbf{Logical}
\(\longrightarrow\)
\textenglish{\ttfamily bool}
\qquad
\textbf{Collection}
\(\longrightarrow\)
\textenglish{\ttfamily list, tuple, set, dict}

\medskip

\textbf{Mutable}
\(\longrightarrow\)
\textenglish{\ttfamily list, set, dict}

\qquad

\textbf{Immutable}
\(\longrightarrow\)
\textenglish{\ttfamily int, float, complex, bool, str, tuple, NoneType}

\end{center}

\end{rememberbox}

\small

%================================================
% 5.35 Input and Output
%================================================

